% 第9章第7节：一般矩阵奇异子空间的扰动（截断处的奇异值间隙，由增广矩阵读出的偏转界，最优误差与最优方向）。

\section{一般矩阵奇异子空间的扰动}
\label{sec:ls-singular-subspace-sensitivity}

上一节说明，奇异值本身总是稳定的。按奇异值选出的左、右奇异子空间却与对称矩阵的特征子空间一样，可能在小扰动下明显转动。本节借助对称增广矩阵，把第~\ref{sec:ls-subspace-sensitivity}节的谱隙界移到奇异子空间上，再对比最优低秩逼近的误差与所用方向的稳定性。

\subsection{截断处的奇异值间隙}
奇异子空间可以通过对称增广矩阵回到特征子空间。设$\bm A,\bm E\in\mathbb R^{m\times n}$，$p=\min(m,n)$，$1\leq k<p$，前$k$个左、右奇异向量排成$\bm U_k$、$\bm V_k$。要求截断处的奇异值间隙$\delta_\sigma=\sigma_k(\bm A)-\sigma_{k+1}(\bm A)>0$，这也保证了所追踪的奇异值$\sigma_1,\ldots,\sigma_k$都为正。若$\sigma_k=0$，引理~\ref{lem:ls-symmetric-augmentation-spectrum}中的$\bm w_k^+$与$\bm w_k^-$对应同一个特征值$0$，就无法从增广矩阵的特征子空间中分出相应的奇异方向。

由引理~\ref{lem:ls-symmetric-augmentation-spectrum}，$\bm H(\bm A)$前$k$个特征值为$\sigma_1,\ldots,\sigma_k$，第$k+1$个为$\sigma_{k+1}$，对应的特征子空间由$\bm w_i^+$（$i\leq k$）张成，其正交投影为
\[
 \bm P_{\bm H}=\frac12\begin{bmatrix}\bm U_k\bm U_k^\top&\bm U_k\bm V_k^\top\\\bm V_k\bm U_k^\top&\bm V_k\bm V_k^\top\end{bmatrix}.
\]

\subsection{由增广矩阵读出的偏转界}
\begin{theorem}[奇异子空间的谱隙界]
\label{thm:ls-singular-subspace-gap}
在上述记号下，设$\delta_\sigma=\sigma_k(\bm A)-\sigma_{k+1}(\bm A)>0$，$\lVert\bm E\rVert_2\leq\varepsilon<\delta_\sigma/2$，$\widetilde{\bm U}_k$、$\widetilde{\bm V}_k$是$\bm A+\bm E$的前$k$个左、右奇异向量。则
\[
 \begin{aligned}
 \lVert\bm U_k\bm U_k^\top-\widetilde{\bm U}_k\widetilde{\bm U}_k^\top\rVert_2&\leq\frac{2\varepsilon}{\delta_\sigma-\varepsilon}\leq\frac{4\varepsilon}{\delta_\sigma},\\
 \lVert\bm V_k\bm V_k^\top-\widetilde{\bm V}_k\widetilde{\bm V}_k^\top\rVert_2&\leq\frac{2\varepsilon}{\delta_\sigma-\varepsilon}\leq\frac{4\varepsilon}{\delta_\sigma}.
 \end{aligned}
\]
\end{theorem}
\begin{proof}
由定理~\ref{thm:ls-weyl-singular-value}，$\bm A+\bm E$在截断处同样分开，对应的投影$\widetilde{\bm P}_{\bm H}$有相同的分块形式。$\bm H(\bm E)$的谱范数等于$\lVert\bm E\rVert_2$，对增广矩阵使用式~\eqref{eq:ls-gap-subspace-budget}，得$\lVert\bm P_{\bm H}-\widetilde{\bm P}_{\bm H}\rVert_2\leq\varepsilon/(\delta_\sigma-\varepsilon)$。矩阵左上块的谱范数不超过整个矩阵的谱范数（对单位向量$\bm x$，左上块乘$\bm x$是整个矩阵乘$(\bm x;\bm0)$所得向量的上半部分），而$\bm P_{\bm H}-\widetilde{\bm P}_{\bm H}$的左上块是$\tfrac12(\bm U_k\bm U_k^\top-\widetilde{\bm U}_k\widetilde{\bm U}_k^\top)$，它的谱范数不超过$\varepsilon/(\delta_\sigma-\varepsilon)$，乘以$2$即得左奇异子空间的结论；右下块是$\tfrac12(\bm V_k\bm V_k^\top-\widetilde{\bm V}_k\widetilde{\bm V}_k^\top)$，同理得右奇异子空间的结论。最后一个不等号用到$\delta_\sigma-\varepsilon>\delta_\sigma/2$。
\end{proof}

可见左、右奇异子空间的稳定性由截断处的间隙$\delta_\sigma$控制，与最小奇异值无关。回到第~\ref{sec:ls-perturbation}节与上一节的$\bm A=\operatorname{diag}(1,s)$，$0<s<1$。取$k=1$，间隙为$\delta_\sigma=1-s$；误差预算$\varepsilon<(1-s)/2$时，第一左、右奇异方向的偏转不超过$4\varepsilon/(1-s)$。$s$很小时矩阵离奇异矩阵很近，这个界反而更好；而$s$接近$1$时，即使矩阵良态，第一奇异方向也容易转动。对半正定矩阵，奇异子空间就是特征子空间，第~\ref{sec:ls-subspace-sensitivity}节示例~\ref{ex:ls-gap-direction-rotation}的转角$\tfrac12\arctan(2t/\delta)$也说明了同样的量级。

\subsection{最优误差与最优方向}
与上一节对照：最佳秩$k$逼近的误差$\sigma_{k+1}$在任何情况下都只移动不超过$\varepsilon$，而取得这个误差所用的主方向只在截断处有足够间隙时才稳定。$\sigma_k=\sigma_{k+1}$时，最优误差仍然稳定，最优的秩$k$矩阵却可能在不同的子空间之间跳变。第~\ref{sec:ls-truncated-svd-ridge}节的截断求解直接用到这个区别：选取截断位置时，既要看被舍弃的奇异值是否可以忽略，也要看保留的方向是否与舍弃的方向分得开。

综合第~\ref{sec:ls-spectral-value-sensitivity}节至本节，把矩阵本身当作输入时：对称矩阵的特征值与任意矩阵的奇异值都随扰动有界地移动，移动量不超过$\lVert\bm E\rVert_2$，并且不要求谱值分开，由此可以可靠地判断符号、正定性和秩的下界；所选特征子空间与奇异子空间的方向要用主夹角比较，偏转由扰动与相应间隙之比控制；一般方阵的特征值与特征方向还取决于非正规的程度，特征值甚至可以不与扰动成比例。条件数与谱的扰动讨论的都是一次作用或一次分解；同一个矩阵的反复作用留给第~\ref{chap:linear-recurrence-stability}章。
